\section{Metodologia}\label{sec:methodology}

A metodologia une a física da atmosfera e a inteligência artificial para lidar com as imagens capturadas pelo INPE, organizada em fases de pesquisa plenamente factíveis de serem executadas no período de 24 meses do Mestrado.

O primeiro passo é o preparo. Devido à iluminação irregular do céu noturno, aplicaremos filtros para suavizar e subtrair o brilho difuso. Em seguida, melhoraremos o contraste para que os contornos das bolhas se destaquem contra o fundo escuro.

Como a maioria das noites não apresenta bolhas, usaremos estratégias para equilibrar os dados, multiplicando imagens e gerando novos exemplos para impedir que o algoritmo fique viciado em prever o céu limpo. Para reduzir o peso dos arquivos sem perder a qualidade espacial da fotografia, aplicaremos a Análise de Componentes Principais Bidimensional, preservando características e acelerando o processamento.

A classificação final será validada por um comitê de inteligência artificial unindo três modelos em equipe. Ferramentas de transparência garantirão que a máquina de fato analisa o formato das bolhas.
